\subsection{Dimension and secant sequences}

Let A be a Noetherian ring, M a finitely generated A-module, S a subset of the radical of A, S the ideal of A generated by S, and a the annihilator of M. We have

\[
\dim_ {\mathrm{A}} (\mathrm{M}) = \dim (\mathrm{A} / \mathrm{a}) \leqslant \dim (\mathrm{A});\tag{6}
\]

Thus, if A is local, we have \(\dim_{\mathbf{A}}(\mathbf{M}) < +\infty\). Moreover, the support of the A-module M/SM is equal to V(a + S) by the corollary of Prop. 18 of II, § 4, no. 4, whence

\[
\dim_ {A} (M / S M) = \dim (A / (a + \mathfrak {S}))  .\tag{7}
\]

When S is finite, or when M is not reduced to 0, we have the inequality

\[
\dim_ {\mathbf {A}} (\mathbf {M} / \mathrm{SM}) \leqslant \dim_ {\mathbf {A}} (\mathbf {M}) \leqslant \operatorname{Card} (\mathbf {S}) + \dim_ {\mathbf {A}} (\mathbf {M} / \mathrm{SM});\tag{8}
\]

When S is finite, this follows from prop. 2 of no. 1 and from formulas (6) and (7) above, and the case where S is infinite is trivial.

DEFINITION 1. --- Let A be a Noetherian local ring, M a nonzero finitely generated A-module, and S a subset of the maximal ideal \(\mathfrak{m}_{\mathrm{A}}\) of A. One says that S is secant for M if one has

\[
\dim_ {A} (M) = \operatorname{Card} (S) + \dim_ {A} (M / S M).\tag{9}
\]

If S is secant for M, one has \(\operatorname{Card}(S) \leqslant \dim_{\mathbf{A}}(\mathbf{M})\) , so S is finite. One says that a family \((x_i)_{i \in I}\) of elements of \(\mathfrak{m}_{\mathbf{A}}\) is secant for M if one has

\[
\dim_ {\mathrm{A}} (\mathrm{M}) = \operatorname{Card} (\mathrm{I}) + \dim_ {\mathrm{A}} \left(\mathrm{M} / \sum_ {i \in \mathrm{I}} x _ {i} \mathrm{M}\right),\tag{10}
\]

that is, if \(i\mapsto x_{i}\) is a bijection from I onto a subset of \(\mathfrak{m}_{\mathbf{A}}\) secant for M.

We say that an element \(x\) of \(m_A\) is secant for \(M\) if \(\{x\}\) is a secant subset for \(M\) , that is, if one has

\[
\dim_ {\mathbf {A}} (\mathbf {M} / x \mathbf {M}) = \dim_ {\mathbf {A}} (\mathbf {M}) - 1.
\]

Remarks. --- 1) It follows from formulas (6) and (7) that S is secant for M if and only if it is secant for A/a, where a is the annihilator of M.

\begin{enumerate}
\def\labelenumi{\arabic{enumi})}
\setcounter{enumi}{1}
\tightlist
\item
  Let S and S′ be two disjoint subsets of \(m_{A}\) . For \(S \cup S'\) to be secant for M, it is necessary and sufficient that S be secant for M and S' be secant for \(M' = M/SM\). This follows from inequality (8) and the formula
\end{enumerate}

\[
\begin{array}{r l} \operatorname{Card} (S \cup S ^ {\prime}) + \dim_ {A} (M / (S M + S ^ {\prime} M)) - \dim_ {A} (M) = & \\ = (\operatorname{Card} (S) + \dim_ {A} (M / S M) - \dim_ {A} (M)) + & \\ + (\operatorname{Card} (S ^ {\prime}) + \dim_ {A} (M ^ {\prime} / S ^ {\prime} M ^ {\prime}) - \dim_ {A} (M ^ {\prime})) \end{array}
\]

\begin{enumerate}
\def\labelenumi{\arabic{enumi})}
\setcounter{enumi}{2}
\tightlist
\item
  Let \(x \in \mathfrak{m}_{\mathbf{A}}\) and \(n \geqslant 1\) be an integer. It is immediate that the modules \(\mathbf{M} / x\mathbf{M}\) and \(\mathbf{M} / x^n\mathbf{M}\) have the same support, hence the same dimension. Consequently, \(x\) is secant for \(\mathbf{M}\) if and only if \(x^n\) is secant for \(\mathbf{M}\). From this and Remark 2, we immediately deduce the following result: let \(x_1, ..., x_r\) be elements of \(\mathfrak{m}_{\mathbf{A}}\) and let \(n_1, ..., n_r\) be integers \(> 0\); then the sequence \((x_1, ..., x_r)\) is secant for \(\mathbf{M}\) if and only if the sequence \((x_1^{n_1}, ..., x_r^{n_r})\) is secant for \(\mathbf{M}\) .
\end{enumerate}

PROPOSITION 3. --- Let A be a Noetherian local ring and M a nonzero finitely generated A-module. For an element x of \(\mathfrak{m}_{\mathrm{A}}\) to be secant for M, it is necessary and sufficient that it belong to none of the minimal elements \(\mathfrak{p}\) of \(\operatorname{Supp}(M)\) such that \(\dim(\mathrm{A}/\mathfrak{p}) = \dim_{\mathrm{A}}(\mathrm{M})\), and it is sufficient that the homothety \(x_{\mathrm{M}}\) of ratio x on M be injective.

Let a be the annihilator of M. To say that x is secant for M means that x is secant for A/a, and the support of M consists of the prime ideals p of A such that \(a \subset p\); if \(x_{M}\) is injective, the image of x in A/a is not a zero divisor in A/a. Prop. 3 then follows from Cor. 2 of Prop. 2 of No. 1 applied to the ring A/a.

COROLLARY. --- Every sequence of elements of \(m_{A}\) which is completely secant for M (A, X, p. 157, def. 2) is secant for M.

Let \((x_{1},\ldots ,x_{r})\) be a sequence of elements of \(\mathfrak{m}_{\mathbf{A}}\) which is completely secant for M. Set \(\mathbf{M}_0 = \mathbf{M}\) and by induction \(\mathbf{M}_i = \mathbf{M}_{i - 1} / x_i\mathbf{M}_{i - 1}\) for \(1\leqslant i\leqslant r\) . By cor. 1 of A, X, p. 160, the homothety with ratio \(x_{i}\) in \(\mathbf{M}_{i - 1}\) is injective, whence \(\dim_{\mathbf{A}}(\mathbf{M}_i) = \dim_{\mathbf{A}}(\mathbf{M}_{i - 1}) - 1\) (for \(1\leqslant i\leqslant r\) ) (prop. 3). We therefore have

\[
\dim_ {\mathbf {A}} (\mathbf {M}) = r + \dim_ {\mathbf {A}} \left(\mathrm{M} / (x _ {1} \mathrm{M} + \dots + x _ {r} \mathrm{M})\right),
\]

hence the sequence \((x_{1},...,x_{r})\) is secant for M.

Remark 4. --- It is not true in general that a secant sequence for M is completely secant for M (p. 87, exerc. 6). We will study later the modules over a Noetherian local ring for which every secant sequence is completely secant.

THEOREM 1. --- Let A be a Noetherian local ring, M a nonzero finitely generated A-module, and S a subset of the maximal ideal \(\mathfrak{m}_{\mathbf{A}}\) of A.

\begin{enumerate}
\def\labelenumi{\alph{enumi})}
\item
  If M/SM is of finite length, we have Card(S) \(\geqslant\) dim \(_{A}\) (M); if S is a secant for M, then Card(S) \(\leqslant\) dim \(_{A}\) (M).
\item
  Every secant subset for M is contained in a maximal secant subset for M.
\item
  The following properties are equivalent:
\end{enumerate}

\begin{enumerate}
\def\labelenumi{(\roman{enumi})}
\item
  S is a maximal secant subset for M;
\item
  S is a secant subset for M and \(\operatorname{Card}(\mathbf{S}) = \dim_{\mathbf{A}}(\mathbf{M})\) ;
\item
  M/SM has finite length and Card(S) = dim\_A(M);
\item
  S is a secant subset for M and M/SM is of finite length.
\end{enumerate}

Since we have \(S \subset m_{A}\) , Nakayama's lemma shows that we have M/SM ≠ \{0\}, whence dim \(_{A}\) (M/SM) ≥ 0 with equality if and only if M/SM has finite length. Assertion \(a\) ) then follows from formulas (8) and (9), as well as the equivalence of properties (ii), (iii), and (iv).

The assertion \(b\) ) follows from the fact that the cardinal of every subset of \(\mathfrak{m}_{\mathbf{A}}\) secant for \(\mathbf{M}\) is bounded above by the integer \(\dim_{\mathbf{A}}(\mathbf{M})\) .

By \(a\) ), every subset secant for M, of cardinal equal to \(\dim_{\mathbf{A}}(\mathbf{M})\) , is maximal. It remains to prove that, if S is secant for M and if \(\operatorname{Card}(\mathbf{S}) < \dim_{\mathbf{A}}(\mathbf{M})\) , then S is not maximal. Let a be the annihilator of M, and B the local noetherian ring A/(a + SA). By cor. 2 of prop. 2 of n° 1, there exists an element \(x\) of \(\mathfrak{m}_{A}\) such that \(\dim(\mathrm{B}/x\mathrm{B}) = \dim(\mathrm{B}) - 1\); whence \(x \notin \mathrm{S}\) . By remark 2, the subset S ∪ \(\{x\}\) of \(\mathfrak{m}_{A}\) is secant for A/a, hence for M by remark 1.

COROLLARY. --- The dimension of M is the smallest of the integers \(d \geqslant 0\) for which there exists a sequence \((x_{1}, ..., x_{d})\) of elements of \(m_{A}\) such that the A-module M/ \(\sum_{i=1}^{d} x_{i}M\) has finite length.

Since \(\varnothing\) is a secant subset for M, Theorem 1, \(b\) ) proves the existence of a maximal secant sequence for M, namely \((x_{1},\ldots,x_{d})\) . But then we have \(d=\dim_{\mathbf{A}}(\mathbf{M})\) and the A-module \(\mathbf{M}/\sum_{i=1}^{d}x_{i}\mathbf{M}\) has finite length by property (iii) of Thm. 1, \(c\) ). Conversely, if \((x_{1}',\ldots,x_{d'}')\) is a sequence of elements of \(\mathfrak{m}_{\mathbf{A}}\) such that the A-module \(\mathbf{M}/\sum_{j=1}^{d'}x_{j}'\mathbf{M}\) is of finite length, we have \(d'\geqslant\dim_{\mathbf{A}}(\mathbf{M})\) by th. 1, \(a\) ).

Recall (III, § 3, n° 2, def. 1) that an ideal q of a Noetherian local ring A is an ideal of definition of A if the q-adic and m\_A-adic topologies of A coincide.

Lemma 2. --- Let A be a Noetherian local ring and q an ideal of A. The following conditions are equivalent:

\begin{enumerate}
\def\labelenumi{(\roman{enumi})}
\item
  q is an ideal of definition of A;
\item
  there exists an integer \(n \geqslant 0\) such that \(\mathfrak{m}_{\mathbf{A}}^{n} \subset \mathfrak{q} \subset \mathfrak{m}_{\mathbf{A}}\) ;
\item
  we have q ≠ A and A/q is an A-module of finite length;
\item
  V(q) is equal to \(\{m_{A}\}\) (in other words, \(m_{A}\) is the only prime ideal of A containing q).
\end{enumerate}

Indeed, the equivalence of (i), (ii), and (iv) was proved in III, § 2, no. 5, and the equivalence of (i) and (iii) follows from IV, § 2, no. 5, cor. 2 to prop. 9.

The corollary of Thm. 1 therefore allows us to state the following scholium:

SCHOLIUM. --- The dimension of a Noetherian local ring A is the smallest integer \(d \geqslant 0\) for which there exists an ideal of definition of A generated by d elements.

